\documentclass[11pt]{article}
\usepackage[margin=1in]{geometry}
\usepackage{graphicx}
\usepackage{amsmath, amssymb}
\usepackage{booktabs}
\usepackage{microtype}
\usepackage{caption}
\usepackage{xcolor}
\usepackage{hyperref}
\hypersetup{colorlinks=true, linkcolor=blue!50!black, urlcolor=blue!50!black}

% Cap any included figure so it fits on the page.
\newcommand{\figfit}[1]{\includegraphics[width=\linewidth,height=0.75\textheight,keepaspectratio]{#1}}

\title{Book 2: Subject-behavioural analysis of the auditory superstitious-perception task}
\author{Auditory SP analyses}
\date{}

\begin{document}
\maketitle
\tableofcontents
\bigskip

%==============================================================================
\section{What this book is and how to read it}
%==============================================================================

Book 1 was about the \emph{stimuli}: how five different acoustic-similarity
metrics describe the relationships between the 300 stimuli used in the
experiment, and how those relationships interact with the experimenter's
target / distractor labelling. The headline result of Book 1 was that the
target / distractor split explains only a tiny fraction of stimulus--stimulus
variance, but it is statistically real (Fisher $F\approx 37.7$, label-axis
variance fraction $0.49\%$, $p<10^{-3}$).

Book 2 is about the \emph{subjects}. Twenty-six subjects across two sites
(Oberlin, $N=16$; UChicago, $N=10$) each saw 300 stimuli (150 in a
\texttt{full\_sentence} block and 150 in an \texttt{imagined\_sentence} block)
and gave a binary ``target / not target'' response per trial, for $7748$
total trials.

The questions for Book 2 are:

\begin{enumerate}
\item \textbf{Alignment.} How well do subjects' responses agree with the
      experimenter's labels, separating sensitivity from response bias?
      (Chapter 16.)
\item \textbf{Internal consistency.} If a subject's responses are not aligned
      with the labels, are they nonetheless \emph{consistent} -- that is,
      do they treat acoustically-similar stimuli the same way?
      (Chapter 17.)
\item \textbf{Predictability from acoustics.} Can the five Book 1 similarity
      metrics predict trial-by-trial responses better than the experimenter
      labels alone? (Chapter 18.)
\item \textbf{Group geometry.} If we build a behavioural distance between
      stimuli purely from how often subjects co-assign them, does the
      resulting geometry look like the label geometry, or like the geometry
      of any acoustic metric? (Chapter 19.)
\item \textbf{Latent strategies.} Do subjects cluster into qualitatively
      different response strategies? (Chapter 20.)
\end{enumerate}

Two pieces of context to keep in mind throughout:

\begin{itemize}
\item \emph{No within-subject stimulus repeats.} A given subject sees each of
      their 300 stimuli exactly once, so classical test--retest reliability is
      not directly available. Internal consistency must therefore be defined
      across stimuli (chapter 17), not across repetitions.
\item \emph{Block-conditional structure.} Each of the two blocks
      (\texttt{full\_sentence}, \texttt{imagined\_sentence}) has its own
      independent pool of 150 stimuli. All similarity-derived features in this
      book are computed within block, so a stimulus is compared only to other
      stimuli a subject could have heard in the same block.
\end{itemize}

\subsection{Trial data schema (the only column-level fact you need)}

The per-subject CSVs have one row per trial with the columns:

\begin{center}
\begin{tabular}{ll}
\toprule
Column & Meaning \\
\midrule
\texttt{Subject Number}              & participant ID \\
\texttt{Block Scheme}                & e.g.\ \texttt{full\_sentence} or \texttt{imagined\_sentence} \\
\texttt{Stimulus Number}             & integer ID of the audio file \\
\texttt{Stimulus Type}               & \textbf{experimenter ground truth} (target / distractor) \\
\texttt{Subject Response}            & subject's binary response (target / distractor) \\
\texttt{Play Count}                  & how many times the audio was replayed \\
\texttt{Trial Start Timestamp}       & onset of the trial \\
\texttt{Play Button Clicked Timestamp} & first replay click \\
\texttt{Audio Start / End Timestamp} & playback boundaries \\
\texttt{Subject Response Timestamp}  & when the response was logged \\
\bottomrule
\end{tabular}
\end{center}

\noindent So \texttt{Stimulus Type} is the \emph{label} (what the experimenter
called it) and \texttt{Subject Response} is the \emph{response} (what the
subject called it). Throughout this book, \texttt{true\_target} is the
binary indicator that \texttt{Stimulus Type == ``target''} and
\texttt{called\_target} is the binary indicator that
\texttt{Subject Response == ``target''}.

%==============================================================================
\section{Chapter 16 -- Signal-detection alignment and response bias}
%==============================================================================

\subsection{Plain-English intent}

A subject can fail to match the experimenter's labels for two very different
reasons. They might genuinely be unable to discriminate targets from
distractors (\emph{low sensitivity}), or they might discriminate well but
prefer one response (\emph{response bias}), e.g.\ calling everything
``target''. Signal-detection theory (SDT) gives us two numbers,
\(d'\) (sensitivity) and \(C\) (criterion), that decouple these.

\subsection{Intuitive formula walkthrough}

Let $H$ be the hit rate (probability of saying ``target'' on a real target)
and $F$ the false-alarm rate (probability of saying ``target'' on a
distractor). Defining $z(\cdot) = \Phi^{-1}(\cdot)$ (the inverse standard
normal),
\[
d' = z(H) - z(F), \qquad C = -\tfrac{1}{2}\bigl(z(H) + z(F)\bigr).
\]
$d'$ measures how separated the underlying ``target'' and ``distractor''
distributions are on whatever internal axis the subject uses, in units of
standard deviations. $C=0$ means a neutral criterion; $C<0$ is liberal
(prefers ``target''); $C>0$ is conservative.

To avoid $\pm\infty$ when $H$ or $F$ hits $0$ or $1$, we use the
log-linear correction (Hautus 1995): replace counts $h, f$ over $n_t, n_d$
trials with
\(\hat H=(h+0.5)/(n_t+1), \hat F=(f+0.5)/(n_d+1)\).
A within-subject permutation null (label permutations) gives a per-subject
$p$-value on $d'$. At the group level we test the across-subject
distribution of $d'$ against zero with a one-sided Wilcoxon signed-rank.

\subsection{Result interpretation}

\begin{figure}[h]
\centering
\figfit{stimuli-analyses/outputs/book2/ch16-sdt/ch16_dprime_forest.png}
\caption{Per-subject $d'$ in each block. Each row is a subject; blue is
\texttt{full\_sentence}, red is \texttt{imagined\_sentence}. Most $d'$
values cluster tightly around zero.}
\end{figure}

The forest plot makes the headline result obvious: \emph{most subjects
performed at or very near chance.} Across all $52$ (subject $\times$ block)
cells the mean accuracy is $0.510$ (range $0.42$--$0.58$), the mean hit rate
is $0.480$ and the mean false-alarm rate is $0.460$.

At the group level there is still a small but reliable signal in the
\texttt{full\_sentence} block: mean $d'=0.091$, Wilcoxon one-sided
$p=0.006$, $t$-test against zero $t=2.61$, $p=0.007$. The
\texttt{imagined\_sentence} block has mean $d'=0.015$ and is statistically
indistinguishable from chance ($p=0.38$). Only $1/26$ subjects in the
imagined block had a permutation-significant $d'$; none did in the full
block individually. The group-level $d'$ in the full block is therefore
real but distributed across all subjects rather than concentrated in a few
strong performers.

\begin{figure}[h]
\centering
\figfit{stimuli-analyses/outputs/book2/ch16-sdt/ch16_dprime_vs_criterion.png}
\caption{Sensitivity ($d'$) versus response bias ($C$). Slight conservative
bias in the imagined block (mean $C=0.094$, $t=2.06$, $p=0.05$); near-neutral
in the full block.}
\end{figure}

\begin{figure}[h]
\centering
\figfit{stimuli-analyses/outputs/book2/ch16-sdt/ch16_roc_per_subject.png}
\caption{Each subject's (FA, hit) point, by block. Points hug the diagonal,
again indicating near-chance discrimination.}
\end{figure}

\subsection{Reader takeaway}

Subjects as a group do barely above chance on the \texttt{full\_sentence}
block, and at chance on the \texttt{imagined\_sentence} block. There is a
small conservative bias in the imagined block. Any further analysis must
therefore explain a behavioural signal that is genuinely small.

%==============================================================================
\section{Chapter 17 -- Internal response consistency}
%==============================================================================

\subsection{Plain-English intent}

If a subject is at $51\%$ accuracy, they could either be (a) randomly
guessing or (b) using a stable but ``wrong'' rule. We cannot distinguish
these from accuracy alone, but we \emph{can} ask whether stimuli that are
acoustically near each other (under some metric) tend to receive the same
response from that subject. If yes, the subject is responding consistently
to acoustic structure even when not matching the labels.

\subsection{Intuitive formula walkthrough}

For each metric $m$ and each block, we build a $k$-nearest-neighbour graph
on the 150 stimuli. Then for each subject $s$ we compute
\[
\textrm{cons}_{s,m} = \frac{1}{N\,k}\sum_{i=1}^{N}\sum_{j\in \mathrm{kNN}(i)}
\mathbf{1}\bigl[Y_{s,i} = Y_{s,j}\bigr],
\]
the average rate at which two acoustically near-neighbour stimuli get the
same response from subject $s$. A within-subject permutation null is built
by shuffling that subject's response vector over the 150 stimuli, which
preserves the marginal response rate. We use $k=10$, $500$ permutations.

\subsection{Result interpretation}

\begin{figure}[h]
\centering
\figfit{stimuli-analyses/outputs/book2/ch17-consistency/ch17_metric_consistency_box.png}
\caption{Distribution across subjects of $k=10$ neighbour consistency, per
metric and block. All medians sit near the chance value of $0.5$.}
\end{figure}

Across both blocks and all five metrics, the per-subject neighbour
consistency clusters tightly around $0.51$ (range of group means
$0.506$--$0.519$ across the ten metric--block cells). This is essentially
indistinguishable from the chance value of $0.50$ that would be expected if
neighbour structure carried no information about the subject's response
choices.

\begin{figure}[h]
\centering
\figfit{stimuli-analyses/outputs/book2/ch17-consistency/ch17_consistency_vs_accuracy.png}
\caption{Per-subject internal neighbour consistency vs.\ label accuracy,
faceted by metric. Quadrants from the planning report: top-right
(consistent and accurate), top-left (consistent but mis-labelled rule),
bottom-right (accurate but inconsistent), bottom-left (chance). Almost
every subject lives near the bottom-left intersection.}
\end{figure}

\subsection{Reader takeaway}

There is no evidence that subjects are using any acoustic-similarity
neighbourhood -- under any of the five metrics from Book 1 -- as the basis
for their responses. Combined with chapter 16, this rules out the
``consistent but unlabelled rule'' explanation: subjects are not just using
a different stable rule based on these acoustic features, they are giving
near-random responses with respect to acoustic neighbourhood.

%==============================================================================
\section{Chapter 18 -- Trial-level logistic prediction}
%==============================================================================

\subsection{Plain-English intent}

The neighbour-consistency analysis is a coarse $k$-NN summary. A logistic
regression is the natural smooth, multi-feature version of the same
question: given each stimulus's relationship to the rest of its block (under
each acoustic metric) and the trial number, can we predict whether this
subject will say ``target''?

\subsection{Intuitive formula walkthrough}

For each metric $m$, define the within-block target--minus--distractor
contrast for stimulus $i$:
\[
x_i^{(m)} = \frac{1}{|T_b\setminus\{i\}|}\sum_{j\in T_b\setminus\{i\}} S^{(m)}_{ij}
            -\frac{1}{|D_b\setminus\{i\}|}\sum_{j\in D_b\setminus\{i\}} S^{(m)}_{ij},
\]
where $T_b, D_b$ are the target / distractor sets of stimulus $i$'s block
and $S^{(m)}$ is the metric-$m$ similarity matrix. Positive $x_i^{(m)}$
means stimulus $i$ is more similar to the block's targets than to its
distractors under metric $m$.

For each (subject, block) we fit
\[
P(Y_t = 1 \mid x_t, \tau_t) = \sigma\!\bigl(\beta_0 + \sum_m \beta_m\, x_t^{(m)} + \beta_\tau \tau_t\bigr),
\]
where $\tau_t$ is the within-block trial index z-scored. We use 5-fold
stratified CV and report the held-out AUC.

\subsection{Result interpretation}

\begin{figure}[h]
\centering
\figfit{stimuli-analyses/outputs/book2/ch18-logistic/ch18_held_out_auc_box.png}
\caption{Per-subject 5-fold held-out AUC for several feature sets, by
block. The $y=0.5$ line is chance.}
\end{figure}

The mean held-out AUCs sit in the very narrow range $0.50$--$0.53$.
Crucially, the \texttt{trial\_order\_only} model -- which uses the within-
block trial index alone -- already achieves $\sim 0.52$ AUC on average,
matching or beating every single-metric model. The
\texttt{true\_label\_baseline} model, which uses only the experimenter's
ground-truth label, gives $0.505$ in the full block and $0.498$ in the
imagined block. The combined \texttt{all\_metrics} model does not help
either ($0.509$ full, $0.520$ imagined).

\begin{figure}[h]
\centering
\figfit{stimuli-analyses/outputs/book2/ch18-logistic/ch18_per_subject_auc_lines.png}
\caption{Paired per-subject AUC: the experimenter true label (left)
vs.\ the all-metrics combined model (right). The all-metrics model is at
best a hair better.}
\end{figure}

\subsection{Reader takeaway}

The five Book 1 acoustic-similarity metrics offer essentially no predictive
purchase on individual responses. The fact that a one-feature
\texttt{trial\_order} model is competitive suggests that what little
structure exists in the responses is closer to a slow drift in
liberal/conservative tendency over a block than a stimulus-locked
acoustic-feature dependence.

%==============================================================================
\section{Chapter 19 -- Response geometry and RSA}
%==============================================================================

\subsection{Plain-English intent}

Even if individual subjects are noisy, perhaps \emph{across} subjects there
is a stable answer to ``which pairs of stimuli get the same response?''. We
build that answer as a behavioural co-assignment matrix and ask which of our
several stimulus geometries (label, Pearson, lagged xcorr, envelope, log
spectrum, MFCC) it most resembles.

\subsection{Intuitive formula walkthrough}

For each block and each pair $(i,j)$ of stimuli,
\[
C_{ij} = \frac{1}{|S_{ij}|}\sum_{s\in S_{ij}} \mathbf{1}\bigl[Y_{s,i}=Y_{s,j}\bigr],
\]
where $S_{ij}$ is the set of subjects who responded on both stimuli. We treat
$1-C$ as a behavioural representational dissimilarity matrix (RDM), then:

\begin{itemize}
\item Cluster on $1-C$ (Agglomerative, average linkage, $k\in\{2,3,4,5\}$)
      and report ARI / NMI vs.\ the experimenter labels.
\item RSA: Spearman correlation of the upper triangle of $1-C$ with the
      upper triangle of each model RDM. Permutation $p$ values come from
      shuffling stimulus indices on the model side ($2000$ permutations).
\end{itemize}

\subsection{Result interpretation}

\begin{figure}[h]
\centering
\figfit{stimuli-analyses/outputs/book2/ch19-rsa/ch19_coassignment_heatmap.png}
\caption{Behavioural co-assignment matrices, with stimuli ordered so that
all targets precede all distractors (white lines mark the boundary). If
behaviour tracked the labels, the diagonal blocks would be visibly lighter
than the off-diagonals; they are not.}
\end{figure}

\begin{figure}[h]
\centering
\figfit{stimuli-analyses/outputs/book2/ch19-rsa/ch19_rsa_bars.png}
\caption{RSA Spearman correlations between the behavioural RDM and each
model RDM, by block, with permutation $p$-values annotated.}
\end{figure}

The Adjusted Rand Index of behaviour-derived clusters against the
experimenter labels is $\le 0.01$ in absolute value for every $k$ in both
blocks. RSA correlations against all six model RDMs (label plus five
acoustic metrics) cluster around $|r|<0.005$ and are not significant under
permutation, with one curious exception: in the
\texttt{imagined\_sentence} block, the behavioural RDM correlates
\emph{negatively} with the log-spectrum RDM ($r=-0.029$, $p=0.006$),
i.e.\ stimuli with similar log spectra are slightly \emph{less} likely to
share responses. The effect is tiny and we flag it only as a hypothesis to
revisit if more data become available.

\subsection{Reader takeaway}

The geometry of subject responses is essentially flat with respect to all
six candidate stimulus geometries. Whatever drives the per-subject responses
is not shared at the group level in any way that respects acoustic
neighbourhood or experimenter labelling.

%==============================================================================
\section{Chapter 20 -- Latent response strategies}
%==============================================================================

\subsection{Plain-English intent}

Maybe subjects \emph{do} use distinct strategies, just at the individual
level. If so, clustering subjects by their full response vector (across all
150 stimuli of a block) might reveal e.g.\ a ``yes-saying'' group and a
``label-tracking'' group. We complement that with clustering on the
metric-weight profiles fit per subject in chapter 18, in case different
subjects rely on different acoustic cues even when their overall accuracy
is the same.

\subsection{Intuitive formula walkthrough}

\textbf{(A) Response-vector clustering.} Build the subject $\times$ stimulus
matrix $Y$, mean-centre rows, compute pairwise cosine similarity, rescale
to $[0,1]$ as an affinity, and run spectral clustering for $k\in\{2,3,4\}$.

\textbf{(B) Coefficient-profile clustering.} For each subject fit a
five-feature logistic regression (one feature per metric, standardised) on
the subject's trials. Z-score the resulting weight vectors across subjects
and run $K$-means with $k\in\{2,3,4\}$.

\subsection{Result interpretation}

\begin{figure}[h]
\centering
\figfit{stimuli-analyses/outputs/book2/ch20-strategies/ch20_subject_response_pca.png}
\caption{Subjects in the first two principal components of their response
vector, coloured by spectral cluster ($k=2$). The clusters reflect mostly a
single axis of slightly liberal vs.\ slightly conservative response rate.}
\end{figure}

\begin{figure}[h]
\centering
\figfit{stimuli-analyses/outputs/book2/ch20-strategies/ch20_cluster_sdt.png}
\caption{ROC points by subject, coloured by the $k=2$ response-vector
cluster. Cluster centroids are very close together; clusters do not
correspond to meaningfully distinct sensitivities.}
\end{figure}

The $k=2$ response clusters in the \texttt{full\_sentence} block split the
26 subjects into groups of $11$ and $15$ subjects with mean response rates
$0.493$ and $0.463$, mean accuracies $0.522$ and $0.508$. In the
\texttt{imagined\_sentence} block the groups are $10$ vs.\ $16$ subjects
with mean response rates $0.466$ vs.\ $0.462$ and mean accuracies $0.498$
vs.\ $0.500$. Translation: the clusters reflect a tiny amount of structure
in subjects' overall response rate but no qualitative strategy difference.

\begin{figure}[h]
\centering
\figfit{stimuli-analyses/outputs/book2/ch20-strategies/ch20_subject_coef_heatmap.png}
\caption{Per-subject metric weights from the logistic regressions (rows
are subjects, columns are metrics), with rows sorted by the $k=2$ cluster
on the weights. The visible vertical structure is not a coherent
``strategy'' but rather noise around zero -- consistent with the chapter
18 finding that all metric features are near-uninformative.}
\end{figure}

\subsection{Reader takeaway}

There is no evidence for distinct, acoustically-defined response strategies
across these 26 subjects. The mild bimodality the clusterer finds is
explained almost entirely by small differences in overall response rate.

%==============================================================================
\section{Synthesis}
%==============================================================================

The five chapters of Book 2 paint a coherent and empirically tight picture:

\begin{itemize}
\item \textbf{Subjects are at chance overall}, with a small but real
      group-level $d'$ of $0.09$ in \texttt{full\_sentence}
      ($p=0.006$) and $d'\approx 0$ in \texttt{imagined\_sentence}.
\item \textbf{Subjects' responses are not consistent with respect to any
      Book 1 acoustic-similarity neighbourhood} -- the $k=10$ neighbour
      consistency sits at $0.51$ for every metric in every block, where
      $0.50$ is chance.
\item \textbf{Logistic models built from the Book 1 metrics do not predict
      subject responses better than chance} (held-out AUC $\le 0.53$ in all
      cases, including a model that combines all five metrics; a model that
      uses only the within-block trial index does just as well).
\item \textbf{The behavioural co-assignment geometry is essentially flat}
      against all six candidate model geometries (Spearman $|r|<0.03$, mostly
      $<0.005$; ARI vs.\ labels $|<0.01|$).
\item \textbf{No distinct response strategies cluster the subjects} --
      attempted clusterings differ only in overall response rate, with no
      meaningful sensitivity gap.
\end{itemize}

A natural reading of these five results together is: \emph{the experiment as
designed produces near-random responses}, and the small label-related
signal in the \texttt{full\_sentence} block is distributed broadly across
subjects rather than carried by a competent subgroup. Because the Book 1
metrics already showed that the label axis explains less than $0.5\%$ of
stimulus-pair variance (chapter 8), this is consistent with the
psychophysical task being above human discrimination threshold for these
stimuli rather than a failure of any single subject or analysis.

\subsection*{Limitations and what would change the picture}

\begin{itemize}
\item \emph{No within-subject repetitions.} Direct test--retest reliability
      is unavailable, so all consistency claims are at the across-stimulus
      level (chapter 17) or across-subject (chapter 19).
\item \emph{Modest sample.} $26$ subjects $\times$ $150$ trials per block
      gives the group test reasonable power for moderate effects (the
      $d'=0.09$ effect is detectable), but small per-subject AUC differences
      of $\sim 0.02$ are at the edge of detectability with $5$-fold CV on
      $150$ trials.
\item \emph{Acoustic feature space.} We tested five metrics from Book 1.
      A negative result against \emph{these} five does not preclude richer
      learned representations (e.g.\ self-supervised speech embeddings)
      revealing structure, and is a natural next analysis if the question
      is pursued further.
\end{itemize}

\end{document}
